% SPDX-License-Identifier: CC-BY-4.0
\documentclass[11pt,letterpaper]{article}

\usepackage[T1]{fontenc}
\usepackage{amsmath,amssymb}
\usepackage{tikz}
\usepackage{pgfplots}
\pgfplotsset{compat=1.18}
\usepackage{algorithm}
\usepackage{algpseudocode}
\usepackage{booktabs}
\usepackage{enumitem}
\usepackage[hidelinks]{hyperref}

\newcommand{\N}{\mathbb{N}}
\newcommand{\Z}{\mathbb{Z}}
\newcommand{\Q}{\mathbb{Q}}
\newcommand{\R}{\mathbb{R}}
\newcommand{\C}{\mathbb{C}}
\newcommand{\Map}[3]{#1 \colon #2 \to #3}
\newcommand{\Mod}[1]{#1\text{-}\mathrm{Mod}}
\newcommand{\GroupHom}{\mathrm{Hom}_{\mathbf{Grp}}}
\newcommand{\RingHom}{\mathrm{Hom}_{\mathbf{Ring}}}

\newtheorem{theorem}{Theorem}
\newtheorem{lemma}[theorem]{Lemma}
\newtheorem{corollary}[theorem]{Corollary}
\newtheorem{definition}[theorem]{Definition}
\newtheorem{remark}[theorem]{Remark}
\newenvironment{proof}{\noindent\textit{Proof.}}{\hfill\(\square\)}

\title{Signed traps and the residue dichotomy}
\author{}
\date{\today}

\begin{document}
\maketitle

\begin{abstract}
The raw map sends \(n\) to \(n/d\) when \(d\mid n\) and to \(qn+c\)
otherwise, with signed integers \(q,c,n\) and a nonzero signed divisor \(d\).
On an orbit no state of which is divisible by \(d\) the divisor never acts, and
the orbit follows the single affine recursion \(n\mapsto qn+c\). This document
classifies that orbit by the multiplier alone: a zero multiplier reaches the
constant offset, a unit multiplier translates or stays fixed, the negative
unit multiplier runs the reflection pair, an expanding multiplier is fixed
exactly on the affine fixed-point equation and otherwise escapes every bound.
It also identifies divisibility with a walk on \(\Z/m\Z\) for \(m=|d|\) and
proves that the orbit is trapped exactly when that walk avoids the zero class.
\end{abstract}

\section{The raw map and trapped orbits}

\begin{definition}\label{Def:Map}
A parameter record consists of integers \(q\), \(c\) and a nonzero integer
\(d\). Its raw step is the map
\[
T(n)=\begin{cases} n/d, & d\mid n,\\ qn+c, & d\nmid n,\end{cases}
\qquad n\in\Z .
\]
The orbit of a seed \(n\) is the sequence \(T^{j}(n)\), \(j\ge0\).
\end{definition}

\begin{definition}\label{Def:Trapped}
A seed \(n\) is \emph{trapped} for the record when \(d\nmid T^{j}(n)\) for
every \(j\ge0\).
\end{definition}

On a trapped orbit the branch \(n\mapsto n/d\) never runs, so all of the
following statements are about the affine recursion alone. The divison branch
is not forgotten: Section~\ref{Sec:Residue} decides when it fires.

\begin{lemma}\label{Lem:Affine}
If \(d\nmid n\), then \(T(n)=qn+c\).
\end{lemma}

\begin{proof}
The divisibility test in Definition~\ref{Def:Map} fails, so the second branch
of the definition applies.
\end{proof}

\begin{lemma}\label{Lem:Preserved}
If \(n\) is trapped, then \(T(n)\) is trapped.
\end{lemma}

\begin{proof}
Let \(j\ge0\). The iterates of \(T(n)\) are the iterates of \(n\) shifted by
one: \(T^{j}(T(n))=T^{j+1}(n)\). Since \(n\) is trapped, \(d\) does not divide
that state.
\end{proof}

\section{The offset sum}

\begin{definition}\label{Def:Sum}
The offset sum of \(j\) affine steps with multiplier \(q\ne1\) is
\[
T_{j}(q)=q^{0}+q^{1}+\dots+q^{j-1},\qquad T_{0}(q)=0 .
\]
At \(q=1\) the same recursion gives \(T_{j}(1)=j\).
\end{definition}

\begin{lemma}[Shifted geometric series]\label{Lem:Series}
For every \(q\in\Z\) and \(j\ge0\),
\[
(q-1)\,T_{j}(q)=q^{j}-1 .
\]
\end{lemma}

\begin{proof}
Induction on \(j\). At \(j=0\) both sides are zero. The induction step adds
\(q^{j}\) to the sum and \(q^{j}(q-1)=q^{j+1}-q^{j}\) to the product.
\end{proof}

\begin{lemma}[Recursion]\label{Lem:Recursion}
For every \(q\in\Z\) and \(j\ge0\), \(T_{j+1}(q)=q\,T_{j}(q)+1\).
\end{lemma}

\begin{proof}
By Definition~\ref{Def:Sum}, \(T_{j+1}(q)=T_{j}(q)+q^{j}\), and
Lemma~\ref{Lem:Series} with \((q-1)T_{j}(q)=q^{j}-1\) gives
\(qT_{j}(q)+1=T_{j}(q)+q^{j}\).
\end{proof}

\begin{lemma}[Two-step extension]\label{Lem:Extension}
For every \(q\in\Z\) and \(j\ge0\),
\[
T_{j+2}(q)=T_{j}(q)+q^{j}+q^{j+1}.
\]
\end{lemma}

\begin{proof}
Apply the recursion twice: \(T_{j+2}(q)=qT_{j+1}(q)+1=q^{2}T_{j}(q)+q+1\), and
the recursion once more gives the stated form.
\end{proof}

\begin{lemma}[Degenerate multipliers]\label{Lem:Degenerate}
For every \(j\ge0\),
\[
T_{j+1}(0)=1,\qquad T_{j}(1)=j,\qquad T_{2j}(-1)=0,\qquad
T_{2j+1}(-1)=1 .
\]
\end{lemma}

\begin{proof}
At \(q=0\) the sum has its first term \(q^{0}=1\) and every later term zero. At
\(q=1\) the sum is a sum of \(j\) ones. At \(q=-1\) the terms alternate between
one and negative one, so an even number of them cancels and an odd number
leaves one.
\end{proof}

\section{The affine itinerary}

\begin{lemma}\label{Lem:Itinerary}
If \(n\) is trapped, then for every \(j\ge0\)
\[
T^{j}(n)=q^{j}n+c\,T_{j}(q).
\]
\end{lemma}

\begin{proof}
Induction on \(j\). At \(j=0\) the formula reads \(n=q^{0}n+cT_{0}(q)=n\). For
the step, the state \(T^{j}(n)\) is not divisible by \(d\) by
Lemma~\ref{Lem:Preserved}, so Lemma~\ref{Lem:Affine} applies and
\(T^{j+1}(n)=q\,T^{j}(n)+c\). Substituting the induction hypothesis and using
Lemma~\ref{Lem:Recursion} gives \(q^{j+1}n+cT_{j+1}(q)\).
\end{proof}

\section{Classification by the multiplier}

\begin{lemma}[Zero multiplier]\label{Lem:ZeroConstant}
Let \(q=0\) and let \(n\) be trapped. Then \(T^{j}(n)=c\) for every \(j\ge1\),
and \(c\) is itself trapped with \(T(c)=c\).
\end{lemma}

\begin{proof}
For \(j\ge1\), Lemma~\ref{Lem:Itinerary} with \(q=0\) and
Lemma~\ref{Lem:Degenerate} give \(T^{j}(n)=0\cdot n+c\cdot1=c\). At \(j=1\)
this says \(T(n)=c\), so \(d\nmid c\) because \(n\) is trapped; hence
Lemma~\ref{Lem:Affine} applies to \(c\) and \(T(c)=0\cdot c+c=c\). The
orbit of \(c\) is therefore constant at a state \(d\) does not divide, so
\(c\) is trapped.
\end{proof}

\begin{lemma}[Unit multiplier]\label{Lem:OneTranslation}
Let \(q=1\) and let \(n\) be trapped. Then \(T^{j}(n)=n+jc\) for every
\(j\ge0\), and the offset sum is the index.
\end{lemma}

\begin{proof}
Insert \(q=1\) and \(T_{j}(1)=j\) into Lemma~\ref{Lem:Itinerary}.
\end{proof}

\begin{lemma}\label{Lem:Archimedean}
For every real, hence every integer, bound there is a natural number above it.
\end{lemma}

\begin{proof}
The natural numbers are unbounded in \(\Z\), and \(\Z\) is unbounded in
\(\R\).
\end{proof}

\begin{theorem}[Unit multiplier with offset]\label{Thm:OneUnbounded}
Let \(q=1\), let \(c\ne0\) and let \(n\) be trapped. Then for every \(R\in\Z\)
there is \(j\ge0\) with \(R<|T^{j}(n)|\).
\end{theorem}

\begin{proof}
By Lemma~\ref{Lem:OneTranslation}, \(T^{j}(n)=n+jc\), so
\(|T^{j}(n)|\ge j|c|-|n|\) and \(|c|\ge1\). Choosing \(j\) above
\(|n|+R\) by Lemma~\ref{Lem:Archimedean} gives
\(|T^{j}(n)|\ge j-|n|>R\).
\end{proof}

\begin{theorem}[Negative unit multiplier]\label{Thm:NegOne}
Let \(q=-1\) and let \(n\) be trapped. Then \(T^{2j}(n)=n\) and
\(T^{2j+1}(n)=c-n\) for every \(j\ge0\). Moreover \(T(n)=c-n\) and
\(T(c-n)=n\), so the pair \(\{n,c-n\}\) is a genuine two-cycle whenever
\(n\ne c/2\).
\end{theorem}

\begin{proof}
Insert \(q=-1\), the parity identities of Lemma~\ref{Lem:Degenerate} and
\((-1)^{2j}=1\), \((-1)^{2j+1}=-1\) into Lemma~\ref{Lem:Itinerary}. The two
displayed values are the first two states of the orbit, and applying
Lemma~\ref{Lem:Affine} at \(n\) and at \(c-n\) (the latter is not divisible by
\(d\) because it is \(T(n)\) and \(n\) is trapped) gives the two-step return.
The two states differ exactly when \(2n\ne c\).
\end{proof}

\begin{lemma}[Fixed point]\label{Lem:Fixed}
If \(d\nmid n\), then \(T(n)=n\) if and only if \((q-1)n+c=0\).
\end{lemma}

\begin{proof}
By Lemma~\ref{Lem:Affine}, \(T(n)=qn+c\), and \(qn+c=n\) is equivalent to
\((q-1)n+c=0\).
\end{proof}

\begin{theorem}[Expanding multiplier]\label{Thm:Expanding}
Let \(|q|\ge2\), let \(n\) be trapped and let \((q-1)n+c\ne0\). Then for every
\(R\in\Z\) there is \(j\ge0\) with \(R<|T^{j}(n)|\).
\end{theorem}

\begin{proof}
Write \(A=(q-1)n+c\ne0\) and \(B=-c\). Lemma~\ref{Lem:Itinerary} and
Lemma~\ref{Lem:Series} give the identity
\[
(q-1)\,T^{j}(n)=q^{j}A+B
\]
for every \(j\), because \((q-1)cT_{j}(q)=c(q^{j}-1)\). Hence
\[
|q-1|\,\bigl|T^{j}(n)\bigr|=\bigl|q^{j}A+B\bigr|\ge|q|^{j}|A|-|B|,
\]
and \(|q|^{j}\ge j+1\) because \(|q|\ge2\) and \(2^{j}\ge j+1\). Since
\(|A|\ge1\) and \(|A|,|B|\) are fixed integers, choosing \(j\) above
\(|B|+|q-1|(R+1)\) by Lemma~\ref{Lem:Archimedean} makes the right side exceed
\(|q-1|R\), which gives \(|T^{j}(n)|>R\).
\end{proof}

\begin{remark}\label{Rem:NoDivisor}
Theorem~\ref{Thm:Expanding} needs the trapping premise only through
Lemma~\ref{Lem:Itinerary}: on an orbit that divides, the branch \(n\mapsto n/d\)
replaces the affine step and the formula fails from that time on.
\end{remark}

\section{The residue walk}\label{Sec:Residue}

\begin{definition}\label{Def:Residue}
Let \(m=|d|\). The residue walk of the record is the map
\[
\rho(r)=qr+c \bmod m,\qquad r\in\Z/m\Z ,
\]
with residues represented in \(\{0,1,\dots,m-1\}\).
\end{definition}

\begin{lemma}[Prefix cast]\label{Lem:PrefixCast}
If no state \(T^{i}(n)\) with \(i<j\) is divisible by \(d\), then
\[
T^{j}(n)\equiv \rho^{j}(n)\pmod m .
\]
\end{lemma}

\begin{proof}
Induction on \(j\). At \(j=0\) both sides are \(n\bmod m\). For the step, the
state \(T^{j}(n)\) is not divisible by \(d\), so Lemma~\ref{Lem:Affine} gives
\(T^{j+1}(n)=qT^{j}(n)+c\), and reducing modulo \(m\) and using the induction
hypothesis gives \(\rho^{j+1}(n)\).
\end{proof}

\begin{corollary}\label{Cor:Cast}
If \(n\) is trapped, then \(T^{j}(n)\equiv\rho^{j}(n)\pmod m\) for every
\(j\ge0\).
\end{corollary}

\begin{proof}
A trapped seed satisfies the hypothesis of Lemma~\ref{Lem:PrefixCast} at every
\(j\).
\end{proof}

\begin{theorem}[Dichotomy]\label{Thm:Dichotomy}
A seed \(n\) is trapped if and only if \(\rho^{j}(n)\ne0\) for every \(j\ge0\).
\end{theorem}

\begin{proof}
If \(n\) is trapped and \(\rho^{j}(n)=0\), then Corollary~\ref{Cor:Cast} gives
\(T^{j}(n)\equiv0\pmod m\), so \(m\) divides \(T^{j}(n)\), hence \(d\) divides
it and \(n\) is not trapped.

Conversely, suppose the residue walk avoids zero and let \(j\) be the least
index with \(d\mid T^{j}(n)\). Every earlier state is not divisible by \(d\), so
Lemma~\ref{Lem:PrefixCast} applies at \(j\) and gives
\(T^{j}(n)\equiv\rho^{j}(n)\pmod m\). The left side is zero modulo \(m\), so
\(\rho^{j}(n)=0\), contradicting avoidance. The index is well defined by the
well-ordering of the naturals.
\end{proof}

\begin{remark}\label{Rem:ResidueMeaning}
The dichotomy converts the infinite divisibility predicate into avoidance of
the origin by the deterministic walk \(\rho\) on a finite set of size \(m\).
It is an equivalence, so both directions transfer, and it quantifies over
every time, not merely over a horizon.
\end{remark}

\section{Controls and scope}

\begin{remark}\label{Rem:NotClassical}
With \(q=2\) and \(c=1\) the affine branch is \(2n+1\), not the classical
\(3n+1\). The orbit of \(9\) with divisor \(3\) is \(9,3,1,3,1,\dots\), and
\(T(9)=3\ne2\cdot9+1\). The division branch fires at the first state, so
\(9\) is not trapped.
\end{remark}

\begin{remark}\label{Rem:Degenerate}
Degenerate cases fall inside the classification. At \(q=0\) and \(c=1\) the
seed \(2\) with divisor \(3\) is trapped and reaches the fixed point \(1\) after
one step; at \(c=0\) every orbit reaches \(0\) and then divides, so that offset
admits no trapped seed. At \(q=1\) and \(c=0\) every trapped state is fixed. At
\(q=-1\) and \(c=0\) the reflection pair degenerates to a fixed point only at
\(n=0\), and that value is divisible by every divisor, so it is not a trapped
instance. At \(d=1\) every state is divisible and no seed is trapped.
\end{remark}

\begin{remark}\label{Rem:Scope}
Nothing is claimed for a seed that divides at some time
(Remark~\ref{Rem:NoDivisor}). The illustrations in
Remarks~\ref{Rem:ResidueMeaning}, \ref{Rem:NotClassical} and
\ref{Rem:Degenerate} are not extra theorems. The first-division return map
with its residue-dependent exponent, the density statements, the signed core
certificate with its boundary bound, the forced prime-power classification of
the independent base, the criticality ledger and every convergence statement
for a classical target belong to companion families. The correspondence is
Table~\ref{Table:SignedTrapsLean}.
\end{remark}

\begin{table}[htbp]
\centering
\begin{tabular}{ll}
\toprule
Statement & Formal name \\
\midrule
Definition~\ref{Def:Map} & \texttt{RawStep} \\
Definition~\ref{Def:Sum} & \texttt{TrapSum} \\
Definition~\ref{Def:Trapped} & \texttt{NeverDivides} \\
Lemma~\ref{Lem:Affine} & \texttt{RawStepAffine} \\
Lemma~\ref{Lem:Preserved} & \texttt{NeverDividesShift} \\
Lemma~\ref{Lem:Itinerary} & \texttt{NeverDividesIterateAffine} \\
Lemma~\ref{Lem:Series} & \texttt{TrapSumMulSubOne} \\
Lemma~\ref{Lem:Recursion} & \texttt{TrapSumRec} \\
Lemma~\ref{Lem:Extension} & \texttt{TrapSumAddTwo} \\
Lemma~\ref{Lem:Degenerate} & \texttt{TrapSumZero}, \texttt{TrapSumOne}, \\
 & \texttt{TrapSumNegOneEven}, \texttt{TrapSumNegOneOdd} \\
Lemma~\ref{Lem:ZeroConstant} & \texttt{TrapZeroConstant}, \texttt{TrapZeroFixed} \\
Lemma~\ref{Lem:OneTranslation} & \texttt{TrapOneTranslation} \\
Lemma~\ref{Lem:Archimedean} & \texttt{TrapExistsNatAbove} \\
Theorem~\ref{Thm:OneUnbounded} & \texttt{TrapOneUnbounded} \\
Theorem~\ref{Thm:NegOne} & \texttt{TrapNegOneEven}, \texttt{TrapNegOneOdd}, \\
 & \texttt{TrapNegOneTwoCycle} \\
Lemma~\ref{Lem:Fixed} & \texttt{TrapFixedIff} \\
Theorem~\ref{Thm:Expanding} & \texttt{TrapExpandingEscape} \\
Definition~\ref{Def:Residue} & \texttt{ResidueStepZ} \\
Lemma~\ref{Lem:PrefixCast} & \texttt{CastIterateOfPrefix} \\
Corollary~\ref{Cor:Cast} & \texttt{NeverDividesCastIterate} \\
Theorem~\ref{Thm:Dichotomy} & \texttt{ResidueAvoidsZeroIffNeverDivides} \\
Remarks~\ref{Rem:NoDivisor}--\ref{Rem:Scope} & non-formal scope summary \\
\bottomrule
\end{tabular}
\caption{Formal correspondence of the statements.}
\label{Table:SignedTrapsLean}
\end{table}

\end{document}
